\subsection{Логические операции}
Логика: \verb|and| (И), \verb|or| (ИЛИ), \verb|not| (НЕ). В Lua ложными считаются только \verb|false| и \verb|nil| (\hyperlink{term:truthiness}{истинность}). Операторы \verb|and| и \verb|or| обладают \textbf{коротким замыканием}: правая часть вычисляется только при необходимости. Сравнения работают для чисел и строк (лексикографически).
\begin{codefile}{examples/chapter_04/lua_logic_ops.lua}
\end{codefile}
\paragraph{Пояснения.}
- Идиома \verb|(cond and A or B)| выбирает \verb|A| при истинном \verb|cond| и \verb|B| иначе.
- \verb|nil| часто означает «нет данных», поэтому проверяйте значение перед доступом к полям.
- Строки сравниваются лексикографически: \verb|"a" < "b"| истинно.
\subsection*{Задачи для самостоятельного решения}
\begin{enumerate}
  \item Составьте условие «\verb|x| в диапазоне \verb|[0,1]|».
  \item Верните \verb|1| если \verb|b==true|, иначе \verb|0| с \verb|and/or|.
  \item Составьте проверку «строка \verb|s| не пустая» (\verb|s~=""|).
  \item Сравните \verb|"a"| и \verb|"b"|: проверка «меньше» (\verb|<|).
  \item Составьте выражение «\verb|x>0| и \verb|y~=0|».
  \item Запишите де Моргана: \verb|not(a and b)| эквивалент \verb|(not a) or (not b)|.
  \item Реализуйте «исключающее ИЛИ» для \verb|p| и \verb|q|.
  \item Преобразуйте число \verb|n| в булево: \verb|n~=0|.
  \item Составьте логическое выражение: \verb|t.enabled == true or t.mode=="auto"|.
  \item Верните «сигнал опасности», если \verb|voltage<3.5| или \verb|temp>80|.
  \item Напишите предикат «индекс чётный»: \verb|i % 2 == 0|.
  \item Составьте выражение, выбирающее максимум из \verb|a| и \verb|b| через \verb|and/or|.
  \item Используя \verb|not|, инвертируйте булеву переменную \verb|on|.
  \item Составьте условие «\verb|t| не \verb|nil| и \verb|t.value| существует».
  \item Напишите логическую функцию «входит ли символ \verb|c| в набор \verb|{"r","g","b"}|».
\end{enumerate}
